% =====================================================================
% COSC 3337 — Homework 1 — ANSWER TEMPLATE for Part 0 and Part 1
%
% Typesetting Part 1 in LaTeX earns 5 bonus points.  Everything you need is
% already defined below, so you never have to type \sigma_{...} yourself.
%
% To use this:  upload it to Overleaf (or run pdflatex locally), put your name
% in the \student command, fill in the answers, and submit the PDF.
% You do not submit this .tex file.
%
% ---------------------------------------------------------------------
% THE FOUR MACROS YOU NEED
%
%   \Select{condition}{relation}        sigma_condition ( relation )
%   \Project{attributes}{relation}      Pi_attributes ( relation )
%   \Group{groupby}{aggregate}{rel}     groupby gamma_aggregate ( rel )
%   \GroupAll{aggregate}{relation}      gamma_aggregate ( relation )      (no grouping)
%
% Plus two for typesetting names:
%
%   \Rel{track}        a relation or attribute name
%   \Val{Folk}         a string literal, in quotes
%
% And the plain symbols:  \Cross  \cup  \cap  -  \wedge  \vee  \neg
%
% EXAMPLE — "the titles of albums released in 2022 or later":
%
%   \[ \Project{\Rel{title}}{\Select{\Rel{release\_year} \geq 2022}{\Rel{album}}} \]
%
% Note that underscores in names must be written \_ , as in \Rel{release\_year}.
% =====================================================================
\documentclass[11pt]{article}

\usepackage[margin=1in]{geometry}
\usepackage{amsmath,amssymb}
\usepackage{booktabs}
\usepackage{enumitem}
\usepackage{fancyhdr}

\newcommand{\Select}[2]{\ensuremath{\sigma_{#1}\left(#2\right)}}
\newcommand{\Project}[2]{\ensuremath{\Pi_{#1}\left(#2\right)}}
\newcommand{\Rename}[2]{\ensuremath{\rho_{#1}\left(#2\right)}}
\newcommand{\Group}[3]{\ensuremath{{}_{#1}\gamma_{#2}\left(#3\right)}}
\newcommand{\GroupAll}[2]{\ensuremath{\gamma_{#1}\left(#2\right)}}
\newcommand{\Cross}{\ensuremath{\times}}
\newcommand{\Rel}[1]{\ensuremath{\mathtt{#1}}}
\newcommand{\Val}[1]{\ensuremath{\text{``#1''}}}

% Use this for the list of resulting tuples under each answer.
\newenvironment{tuples}{\par\medskip\noindent\textit{Resulting tuples:}\begin{itemize}[nosep, leftmargin=1.6em]}{\end{itemize}\medskip}

\newcommand{\student}{YOUR NAME HERE}

\pagestyle{fancy}
\fancyhf{}
\lhead{\small COSC 3337 — HW 1 — \student}
\rhead{\small Page \thepage}
\renewcommand{\headrulewidth}{0.4pt}

\setlength{\parindent}{0pt}
\setlength{\parskip}{6pt}

% =====================================================================
\begin{document}

\begin{center}
  {\Large\bfseries COSC 3337 — Homework 1}\\[2pt]
  {\large \student}
\end{center}

% =====================================================================
\section*{Part 1 — Relational algebra}

\subsection*{RA1}

\[
  % your expression here, for example:
  % \Project{\Rel{a_j}}{\Select{\Rel{a_i} < 999}{\Rel{r}}}
\]

\begin{tuples}
  \item ...
\end{tuples}

\subsection*{RA2}

\[
  % your expression here
\]

Your explanation here: how many tuples, how many rows in \Rel{artist}, and why
they differ.

\subsection*{RA3}

\[
  % your expression here
\]

\begin{tuples}
  \item ...
\end{tuples}

\subsection*{RA4}

\[
  % your expression here
\]

\begin{tuples}
  \item ...
\end{tuples}

\subsection*{RA5}

\[
  % your expression here.  Grouped aggregation looks like:
    % \Group{\Rel{a_i}}{\text{sum}(\Rel{a_j})}{\Rel{r}}
\]

\begin{tuples}
  \item ...
\end{tuples}

\subsection*{RA6}

\textbf{(a)} Number of tuples, and how you got it:

\textbf{(b)} What is wrong with the expression:

\textbf{(c)} A correct expression:
\[
  % your expression here
\]

\begin{tuples}
  \item ...
\end{tuples}

\subsection*{RA7}

\textbf{(a)} Equivalent? Justification:

\textbf{(b)} Equivalent? Justification, and the general rule:

\end{document}
